% !TEX program = xelatex
% 공통수학2 · Ⅲ. 함수와 그래프 · 1. 함수 · 03 역함수 · 1/2차시
% 학생용 활동지 (답안 없음)
\documentclass[11pt,a4paper]{article}
\usepackage{kotex}
\usepackage{iftex}
\ifPDFTeX\else
  \setmainhangulfont{Noto Serif CJK KR}
  \setsanshangulfont{Noto Sans CJK KR}
\fi
\usepackage[margin=2.1cm]{geometry}
\usepackage{parskip}
\usepackage{enumitem}
\usepackage{array}
\usepackage{tabularx}
\usepackage{xcolor}
\usepackage{amssymb}
\usepackage{tikz}
\usepackage[most]{tcolorbox}
\usepackage{fancyhdr}
\usepackage{titlesec}

\definecolor{goldc}{HTML}{B07C22}
\definecolor{goldstrong}{HTML}{8A5F16}
\definecolor{indigoc}{HTML}{2B3B57}
\definecolor{indigosoft}{HTML}{6E82A6}
\definecolor{linec}{HTML}{D6DACB}
\definecolor{surface2}{HTML}{E4E8DC}
\definecolor{writeline}{HTML}{C7CCBA}
\definecolor{inksoft}{HTML}{52604F}

\titleformat{\section}{\normalfont\sffamily\bfseries\large\color{indigoc}}{\thesection}{0.6em}{}
\titlespacing{\section}{0pt}{14pt}{6pt}
\newtcolorbox{goalbox}{enhanced, breakable, colback=white, colframe=goldc, boxrule=0.5pt, leftrule=3pt, arc=2pt, left=10pt, right=10pt, top=8pt, bottom=8pt}
\newtcolorbox{sheetbox}{enhanced, breakable, colback=white, colframe=linec, boxrule=0.6pt, arc=3pt, left=11pt, right=11pt, top=9pt, bottom=9pt}
\newcommand{\writespace}[1]{%
  \par\nobreak\vspace{3pt}
  \foreach \n in {1,...,#1} {%
    \noindent\textcolor{writeline}{\rule{\linewidth}{0.4pt}}\par\vspace{0.62cm}%
  }
}
\newcommand{\fillblank}[1]{\makebox[#1]{\hrulefill}}
\pagestyle{fancy}\fancyhf{}\renewcommand{\headrulewidth}{0pt}
\fancyfoot[L]{\footnotesize\color{inksoft} 공통수학2 · 2026학년도 2학기 · 지평선고등학교}
\fancyfoot[R]{\footnotesize\color{inksoft} 다음 시간 --- 03 역함수 2/2(그래프)}

\newcommand{\revmap}[1]{%
  \begin{tikzpicture}[scale=0.6, baseline]
    \draw[thick, indigosoft] (0,0) ellipse (0.9 and 1.6);
    \draw[thick, indigosoft] (2.6,0) ellipse (0.9 and 1.6);
    \node[circle, fill=indigoc, inner sep=1.4pt, label=left:{\footnotesize $1$}] (x1) at (0,1) {};
    \node[circle, fill=indigoc, inner sep=1.4pt, label=left:{\footnotesize $2$}] (x2) at (0,0) {};
    \node[circle, fill=indigoc, inner sep=1.4pt, label=left:{\footnotesize $3$}] (x3) at (0,-1) {};
    \node[circle, fill=goldstrong, inner sep=1.4pt, label=right:{\footnotesize $1$}] (y1) at (2.6,1) {};
    \node[circle, fill=goldstrong, inner sep=1.4pt, label=right:{\footnotesize $4$}] (y2) at (2.6,0) {};
    \node[circle, fill=goldstrong, inner sep=1.4pt, label=right:{\footnotesize $7$}] (y3) at (2.6,-1) {};
    \node[above] at (0,1.75) {\footnotesize $X$};
    \node[above] at (2.6,1.75) {\footnotesize $Y$};
    #1
  \end{tikzpicture}}

\begin{document}

\noindent
\begin{minipage}[b]{0.62\linewidth}
  {\sffamily\footnotesize\color{indigosoft} 공통수학2 · Ⅲ. 함수와 그래프 · 1. 함수}\\[2pt]
  {\sffamily\bfseries\Large 03 역함수 \textperiodcentered\ 36차시}
\end{minipage}\hfill
\begin{minipage}[b]{0.36\linewidth}
  \raggedleft\small\color{inksoft}
  반\ \fillblank{1.6cm} \quad 번호\ \fillblank{1.6cm}\\[4pt]
  이름\ \fillblank{3.6cm}
\end{minipage}

\vspace{4pt}
\noindent{\color{goldc}\rule{\linewidth}{2.2pt}}
\vspace{10pt}

\begin{goalbox}
\textbf{오늘의 목표}\\
역함수의 존재 조건을 이해하고, 함수의 역함수를 구하며 역함수의 성질을 설명할 수 있다.
\vspace{4pt}
\small\color{inksoft}
$\square$ 자신 있음 \quad $\square$ 조금 헷갈림 \quad $\square$ 처음 봄 \hfill (수업 전 나의 상태에 표시)
\end{goalbox}

\section{생각 열기 --- 화살표를 거꾸로}

\begin{sheetbox}
아래 대응에서 화살표의 방향을 모두 반대로 뒤집어 $Y\to X$ 대응을 만들어 보자.

\vspace{4pt}
\begin{center}
\revmap{
  \draw[->, thick] (x1) -- (y1);
  \draw[->, thick] (x2) -- (y2);
  \draw[->, thick] (x3) -- (y3);
}
\end{center}

뒤집은 대응도 여전히 함수(각 원소에 오직 하나씩 대응)라고 할 수 있을까? 원래 대응이 어떤 성질을 가지고 있었기에 그런지 써 보자.
\writespace{3}
\end{sheetbox}

\section{개념 정리 --- 역함수의 뜻}

\begin{sheetbox}
함수 $f:X\to Y$가 \fillblank{2.6cm}일 때, $Y$의 각 원소 $y$에 대하여 $f(x)=y$를 만족하는 $X$의 원소 $x$를 대응시키면 이것도 함수가 된다. 이 함수를 $f$의 \fillblank{2cm}라 하고 $f^{-1}$로 나타낸다.

\vspace{6pt}
\[y=f(x) \iff x=f^{-1}(\ \fillblank{1.4cm}\ )\]

\vspace{6pt}
\textbf{역함수를 구하는 절차}
\begin{enumerate}[leftmargin=1.8em, itemsep=2pt]
  \item $y=f(x)$를 \fillblank{2.6cm} 꼴로 정리한다.
  \item 식에서 $x$와 $y$를 \fillblank{2.4cm}.
\end{enumerate}

\vspace{8pt}
\textbf{예제} --- $f(x)=3x-2$의 역함수\\
$y=3x-2 \;\Rightarrow\; x=\fillblank{2.4cm} \;\Rightarrow\; f^{-1}(x)=\fillblank{2.4cm}$

\vspace{8pt}
\textbf{역함수의 성질}\\
$(f^{-1})^{-1}=\fillblank{1.4cm}$, \quad $f^{-1}\circ f = \fillblank{1.6cm}$(항등함수), \quad $f\circ f^{-1}=\fillblank{1.6cm}$(항등함수)\\
$(g\circ f)^{-1} = \fillblank{2.6cm}$ (순서가 바뀜에 주의)
\end{sheetbox}

\section{연습 문제}

\begin{sheetbox}
\textbf{1.} $f(x)=x^2$ ($X=\mathbb{R}$)이 역함수를 가질 수 없는 이유를 반례를 들어 설명하시오.
\writespace{3}

\vspace{4pt}
\textbf{2.} $f(x)=3x-2$의 역함수 $f^{-1}(x)$를 구하시오.
\writespace{3}

\vspace{4pt}
\textbf{3.} $f(x)=2x+1$일 때, $f^{-1}(5)$와 $f^{-1}(f(3))$을 각각 구하시오.
\writespace{4}

\vspace{4pt}
\textbf{4.} $f(x)=x+1$, $g(x)=2x$일 때, $(g\circ f)^{-1}(x)$를 ⑴ 직접 구하는 방법과 ⑵ $f^{-1}\circ g^{-1}$을 이용하는 방법으로 각각 구하여 비교하시오.
\writespace{5}
\end{sheetbox}

\section{사고의 가시화 --- See · Think · Wonder}

\noindent{\footnotesize\color{inksoft} 오늘 배운 `역함수의 뜻과 성질'을 떠올리며 아래 세 칸을 채워 보자.}\\[6pt]
\begin{tabularx}{\linewidth}{@{}XXX@{}}
\textbf{\footnotesize\color{indigoc} See --- 무엇을 보았나?} &
\textbf{\footnotesize\color{indigoc} Think --- 무엇을 생각했나?} &
\textbf{\footnotesize\color{indigoc} Wonder --- 무엇이 궁금한가?} \\[4pt]
\end{tabularx}
\noindent
\begin{minipage}[t]{0.32\linewidth}\writespace{3}\end{minipage}\hfill
\begin{minipage}[t]{0.32\linewidth}\writespace{3}\end{minipage}\hfill
\begin{minipage}[t]{0.32\linewidth}\writespace{3}\end{minipage}

\section{마무리 성찰}

\begin{sheetbox}
\begin{minipage}[t]{0.48\linewidth}
{\footnotesize\color{inksoft} 역함수가 존재하기 위한 조건을 한 문장으로}
\writespace{2}
\end{minipage}\hfill
\begin{minipage}[t]{0.48\linewidth}
{\footnotesize\color{inksoft} 감사 일기 / 유념 한 줄}
\writespace{2}
\end{minipage}
\end{sheetbox}

\end{document}
